
\IN{full transformation semigroup}
\IN{domain!of transformation semigroup}
\IN{transformation semigroup} 
\IN{transformation monoid}
\label{def:T_X}
\label{def:T_n}
The set of all transformations on a set~$X$
with function composition as operation forms
a semigroup $(X^X, \circ)$, called
the \emph{full transformation semigroup on~$X$}.
This set also forms a monoid with identity $\operatorname{id}_X$.
We also denote $(X^X, \circ)$ by $\mathcal T_X$.
For $X = [n]$, we set $\mathcal T_n = \mathcal T_X$. 
A \emph{transformation semigroup} (resp.~\emph{transformation monoid}) on the \emph{domain}~$X$ 
is a subsemigroup (resp.~submonoid) of $\mathcal T_X$.

\IN{fixpoint!of transformation semigroup}
\IN{$S$-congruence}
\IN{congruence class}
\IN{trivial congruences}
\IN{congruent elements}
Let $S \le_S \mathcal T_X$ be a transformation semigroup. 
A \emph{fixpoint} of $S$ is a point fixed by all elements from $S$. 
An \emph{$S$-congruence} is an equivalence relation~$\sim$ on~$X$ such that,
for all $x,y \in X$ and $t \in S$, we have 
\begin{equation}
	\label{prelim:trans_semigrp_congruence} 
	x \sim y \Rightarrow \applyfunc{x}{t} \sim \applyfunc{y}{t}.
\end{equation} 
We also say that $x$ and $y$ are \emph{congruent} if $x \sim y$. Equivalence classes of~$\sim$
are also called \emph{congruence classes} in this case.
The universal and equality relations on $X$ are $S$-congruences, called the \emph{trivial $S$-congruences}. 

Recall from \cref{prelim:sec:fun} that the elements of $X$ are also called \emph{points} in the context
of transformations.

\begin{remark} \label{prelim:rem:monoid-to-transfmonoid}
	Let $(M, \cdot, 1)$ be a monoid.
	For $m \in M$ let $\hat{m} \colon M \to M$
	be given by $\hat{m}(x) = xm$ for $x \in M$.
	Then $\hat{M} = \{ \hat{m} \mid m \in M \}$ is a transformation monoid
	on the domain~$M$.
	The map  $m \mapsto \hat{m}$
	is injective
	since $\hat{m} = \hat{n}$ for $m,n \in M$ 
	implies $m = \hat{m}(1) = \hat{n}(1) = n$,
	and it is a monoid morphism.
	Thus, $\hat{M}$ represents the monoid~$M$
	as a transformation monoid on the domain~$M$.
	The $M$-congruences of $\hat{M}$ are also called
	right-congruences of~$M$ in the literature, see, for example, \cite{howie1995fundamentals}.
\end{remark}

\begin{remark} \label{prelim:rem:semigrp-to-transfsemigrp}
	Let $S$ be a semigroup.
	By adjoining an identity as in \cref{prelim:rem:semigroup-adjoin-identity} 	
	we obtain a monoid $M = S \cup \{1\}$.
	Consider the transformation monoid
	constructed from it as in \cref{prelim:rem:monoid-to-transfmonoid}
	via the map $m \mapsto \hat{m}$ for $m \in M$.
	Then $\{ \hat{s} \mid s \in S\}$
	is a transformation semigroup and the map $s \mapsto \hat{s}$ for $s \in S$
	is an injective semigroup morphism.
	Thus, the semigroup~$S$ is represented as a transformation semigroup.
\end{remark}

The embedding of an arbitrary semigroup 
in a transformation semigroup as in \cref{prelim:rem:semigrp-to-transfsemigrp}
is also known as Cayley's theorem, 
for example, see~\cite[Thm.~2.4.3]{GanyushkinMazorchuk2008}.
However, Cayley only gave it for groups~\cite{Cayley54a}.\footnote{Cayley did not give an 
	explicit statement, but only implicitly assumed it, since the notion of an abstract group
	was not yet invented. He investigated a formal calculus of operations with an identity element and satisfying associativity (implicitly assuming inverses) and used it to investigate
	groups.}
The extension to semigroups is due to Suschkewitsch~\cite{Suschkewitsch1926}.


